\section{Empirical Results}
\label{sec:feasibility}

We report two results on a \emph{single frozen} speech model, one positive and one negative, and read both through a single quantitative lens---the \emph{spread} of the verifiable reward across candidates. All ``training-free RL'' here is inference-time \emph{verifiable-reward selection}: a frozen bi-encoder (\texttt{omni-embed-nemotron-3b}, a SentenceTransformer) ranks candidate cards by cosine similarity under a verifiable reward, realising best-of-$N$ / reranking / candidate-card retrieval with \emph{no weight updates and no structural change}. The positive result is a content/intent \emph{selection} gain where the reward has large spread across candidates (\Cref{sec:res:content}); the negative is a paralinguistic read-out and a null emotion-selection gain, where the reward has little spread to exploit (\Cref{sec:res:para}). The two are explained by the same lens, which we state first (\Cref{sec:res:lens}).

\subsection{The Reward-Spread Lens}
\label{sec:res:lens}

The optimization we perform is the KL-regularised tilt of the frozen base distribution $\qo$ toward a verifiable reward $\Rwd$: maximising $\Fobj(q)=\Exp_{q}[\Rwd]-\beta\,\KL(q\,\|\,\qo)$ has optimum $\qs \propto \qo\,\exp(\Rwd/\beta)$, the same object that underwrites best-of-$N$ and reward-guided decoding \citep{korbak2022kl,mudgal2024controlled,beirami2024bon}, here applied to selection among candidate cards. The gain of this tilt admits the identity
\[
  \gain \;=\; \beta\,\KL(\qo\,\|\,\qs),
\]
which we use \emph{qualitatively}: it is the Donsker--Varadhan / free-energy identity for the log-partition function $\log\Zpart=\log\Exp_{\qo}[\exp(\Rwd/\beta)]$, not a novel theorem of ours. Its content for us is a comparison, not a constant: the gain is controlled by the reward \emph{spread} $\spread=\max\Rwd-\min\Rwd$ over $\mathrm{supp}(\qo)$. A flat reward---one that does not discriminate among candidates---yields \emph{zero} gain; a reward with genuine spread across candidates yields \emph{strictly positive} gain, and does so \emph{on a single frozen model}, with no decomposition or additional structure required. Two strict, \texttt{sorry}-free Lean lemmas make exactly these two endpoints precise: \texttt{flat\_no\_gain} (a constant reward gives $\gain=0$) and \texttt{gain\_pos\_of\_nonconstant} (a non-constant reward gives $\gain>0$). We claim nothing beyond these two facts; they suffice to organise the evidence. Content and intent, we will see, present the selector with high reward spread and yield selection gains; the paralinguistic read-out of this frozen encoder presents little spread and yields none.

\subsection{Content/Intent Selection: High Spread, Real Gains}
\label{sec:res:content}

Our positive result is a frozen-bi-encoder \emph{selection} gain on spoken-language-understanding intent. The mechanism is fixed and simple: the frozen \texttt{omni-embed-nemotron-3b} encoder embeds the audio query and a set of candidate \emph{cards} (textual schemas describing each tool/intent), and we select the card of maximum cosine similarity under the verifiable intent-match reward. Nothing is decoded and nothing is trained---the ``policy'' is the choice of candidate-card representation the selector ranks. We stress this attribution because it is the whole of the claim: these are \emph{frozen bi-encoder verifiable-selection} gains over candidate cards, \emph{not} generative candidates, best-of-$N$ decoding, or any generative operator, none of which we ran.

On \textbf{MInDS-14} (en-US tool/intent, $n{=}182$, seed $42$), \Cref{tab:content} reports three candidate-card representations under the identical selection rule. A naive card reaches $\mathrm{Acc@1}=0.720$; a raw-schema card raises this to $0.857$; and the policy card---a boundary-aware schema (name, description, examples, boundary notes)---reaches $0.984$. The policy improvement over the raw-schema baseline is $+0.126$, with a $95\%$ paired-bootstrap CI of $[0.077,0.181]$ ($n{=}182$, seed $42$). The gain is well separated from zero: the selector is exploiting real reward spread across the candidate cards. \textbf{SLURP intent} corroborates the same mechanism on a second corpus: the identical frozen-bi-encoder card-selection move improves intent accuracy from $0.550$ to $0.880$ ($+0.330$). Both numbers are single-frozen-model selection gains and are attributed as such.

\begin{table}[t]
\centering
\caption{Frozen bi-encoder verifiable-reward \emph{selection} on intent (single frozen \texttt{omni-embed-nemotron-3b}, cosine ranking of candidate cards; seed $42$). MInDS-14 en-US tool/intent, $n{=}182$: the policy card improves over the raw-schema baseline by $+0.126$, $95\%$ paired-bootstrap CI $[0.077,0.181]$ (committed artifact \texttt{\_repro/minds14\_toolintent\_paired.json}). SLURP intent corroborates the same card-selection mechanism. These are selection gains, not generative decoding.}
\label{tab:content}
\begin{tabular}{llccc}
\toprule
Task ($n$) & Candidate-card representation & Acc@1 & $\Delta$ vs.\ baseline & $95\%$ CI \\
\midrule
MInDS-14 ($182$) & naive card                 & 0.720 & ---     & --- \\
                 & raw-schema (baseline)      & 0.857 & ---     & --- \\
                 & policy (boundary-aware)    & 0.984 & $+0.126$ & $[0.077,0.181]$ \\
\midrule
SLURP intent     & basic-label (baseline)     & 0.550 & ---     & --- \\
                 & boundary card (policy)     & 0.880 & $+0.330$ & --- \\
\bottomrule
\end{tabular}
\end{table}

Reproduce with the committed command against the committed artifact \texttt{\_repro/minds14\_toolintent\_paired.json}. We add an honest contamination caveat: SLURP and MInDS-14 are public benchmarks that may overlap pretraining corpora, so the absolute accuracies should be read with that risk in mind; the \emph{contrast} between candidate-card representations under a fixed selector is the load-bearing quantity, and it is a within-model comparison.

We differentiate this result honestly from adjacent best-of-$N$/reranking work in this series, which targets ASR. The novelty here is narrow and specific: a \emph{frozen omni-embedding} encoder used for SLU/intent \emph{selection} over candidate cards, not ASR best-of-$N$ decoding. The value is the demonstration that a high-spread verifiable reward produces a real, reproducible selection gain on a single frozen speech encoder.

\subsection{Paralinguistics: Low Spread, No Gain}
\label{sec:res:para}

The negative result is equally concrete and, we argue, equally a contribution: a precise committed measurement plus the spread explanation for it. Two facts.

First, the frozen omni \emph{content} read-out carries little paralinguistic information. On CREMA-D, with the read-out layer selected on dev and three seeds, speaker identity is recoverable at \emph{near chance}---mean accuracy $0.033$ on a $91$-way task ($3.0\times$ chance, chance $0.011$)---and emotion only modestly---mean accuracy $0.405$ on a $6$-way task ($2.4\times$ chance, chance $0.167$) (committed artifact \texttt{\_repro/paralinguistic\_negative\_probe.json}). \Cref{tab:para} summarises. For a \emph{content} embedding, near-chance speaker recovery is close to definitional rather than surprising; the point is that an index built on this frozen omni content embedding would retrieve on content, not on who-spoke or how-they-felt.

\begin{table}[t]
\centering
\caption{Frozen omni content read-out on CREMA-D (linear/kNN probe, dev-selected layer, $3$ seeds; committed artifact \texttt{\_repro/paralinguistic\_negative\_probe.json}). The content embedding carries near-chance speaker and only modest emotion information: little paralinguistic spread for a selector to exploit.}
\label{tab:para}
\begin{tabular}{lcccc}
\toprule
Factor & Ways & Chance & Mean acc & $\times$ chance \\
\midrule
Speaker & $91$ & $0.011$ & $0.033$ & $3.0\times$ \\
Emotion & $6$  & $0.167$ & $0.405$ & $2.4\times$ \\
\bottomrule
\end{tabular}
\end{table}

Second, training-free selection over this encoder yields \emph{no significant} paralinguistic gain. Treating the choice of read-out (layer and pooling) as the selectable degree of freedom, and evaluating the emotion improvement of attentive-statistics pooling over a mean-pool baseline---both layers selected on dev and frozen before test, with a paired-delta bootstrap and five seeds $\{42,7,123,2024,31337\}$---the across-seed mean improvement is $\Delta=+0.037$ with a $95\%$ $t$-CI of $[-0.043,+0.116]$, which \emph{spans zero} (committed artifact \texttt{\_repro/emotion\_pool\_paired\_v2.json}). \Cref{tab:emotion-multiseed} reports the per-seed detail. The headline is a null: training-free pool-selection produces no significant emotion gain. We note plainly that an earlier single-seed report of $+0.097$ was the best seed under \emph{oracle test-layer} selection (a selection leak); under dev-selection it reproduces only as the best of five seeds, and the appropriate across-seed aggregate is the null above.

\begin{table}[t]
\centering
\caption{CREMA-D emotion, training-free pool-selection (mean-pool baseline vs.\ attentive-statistics pooling, both layers dev-selected, paired-delta bootstrap on TEST, five seeds; committed artifact \texttt{\_repro/emotion\_pool\_paired\_v2.json}). The across-seed $95\%$ $t$-CI $[-0.043,+0.116]$ spans zero: no significant emotion gain. The earlier single-seed $+0.097$ was the best seed under oracle test-layer selection, reproduced here only as the best of five.}
\label{tab:emotion-multiseed}
\begin{tabular}{lllccrcl}
\toprule
Seed & base (dev-sel) & sel (dev-sel) & acc$_\text{base}$ & acc$_\text{sel}$ & $\Delta$ & paired $95\%$ CI & sig. \\
\midrule
$42$    & audio-mean L0 & audio\_attn L16 & 0.400 & 0.497 & $+0.097$ & $[+0.027,+0.167]$ & SIG \\
$7$     & audio-mean L0 & audio\_attn L0  & 0.423 & 0.370 & $-0.053$ & $[-0.117,+0.010]$ & n.s. \\
$123$   & audio-mean L0 & audio\_attn L16 & 0.390 & 0.473 & $+0.083$ & $[+0.013,+0.150]$ & SIG \\
$2024$  & audio-mean L0 & audio\_attn L16 & 0.493 & 0.487 & $-0.007$ & $[-0.077,+0.063]$ & n.s. \\
$31337$ & audio-mean L0 & audio\_attn L16 & 0.400 & 0.463 & $+0.063$ & $[-0.010,+0.130]$ & n.s. \\
\midrule
mean    & --- & --- & --- & --- & $+0.037$ & across-seed $95\%$ $t$-CI $[-0.043,+0.116]$ (spans $0$; null) & --- \\
\bottomrule
\end{tabular}
\end{table}

We keep the scope of this negative tight. It is a linear/kNN probe and a pool-selection experiment on \emph{one} frozen content bi-encoder; it does not generalise to frozen omni models as a class, and CREMA-D content is a $12$-sentence closed read set, so we make no content claim from it. What it establishes is precisely what the spread lens predicts: this content encoder presents the selector with little paralinguistic spread, so training-free selection over it has little to exploit.

\subsection{Reading the Evidence}
\label{sec:res:reading}

The two results are one story. The gain of training-free selection tracks the reward \emph{spread} across candidates, exactly as the OSA-1 identity $\gain=\beta\,\KL(\qo\,\|\,\qs)$ and its two Lean endpoints (\texttt{flat\_no\_gain}, \texttt{gain\_pos\_of\_nonconstant}) describe. Content and intent give the selector high spread---the candidate cards differ substantially in verifiable reward---and we obtain real, reproducible single-model gains ($+0.126$ $[0.077,0.181]$ on MInDS-14; $+0.330$ corroborating on SLURP). The paralinguistic read-out of the same frozen encoder gives the selector little spread, and we obtain no significant gain (emotion across-seed CI $[-0.043,+0.116]$; speaker near chance). We report both intervals in full and state the negative as plainly as the positive.

A brief note on what an \emph{agent-level} extension would require, to be honest about scope rather than to claim it. Moving beyond single-model selection to a genuinely compositional, cross-session program would need (a) a cross-session same-speaker corpus, which our frozen data does not contain; (b) a locally-driven generative operator, which we never ran; and (c) a genuinely new non-separable irreducibility result. Absent these, the spread lens is already informative: on a fixed reward, mere context isolation adds no spread and therefore no gain. We leave that program to future work and rest this paper on the two measured results above.